% ---
% Inicia os anexos
% ---
\begin{anexosenv}
    \chapter{Código Exercício 2}\label{cap:codigo_ex2}
    
    \begin{lstlisting}[language=Python, breaklines=true, basicstyle=\tiny]
    import numpy as np
    import matplotlib.pyplot as plt
    import pandas as pd

    # Parameters
    e0 = 8.854187817e-12  # Vacuum permittivity (F/m)
    mu0 = 4 * np.pi * 1e-7  # Vacuum permeability (H/m)
    vp = 1 / np.sqrt(e0 * mu0)

    a = 1 # tamanho da reta (eixo x)
    T = 10e-9 # tempo de simulação 

    dx = 0.01
    dt = 0.005e-9

    nx = int(a / dx) + 1
    nt = int(T / dt)

    xgrid = np.linspace(0, a, nx)

    # campos
    Ey = np.zeros((nt, nx))
    Hz = np.zeros((nt, nx))

    # valor inicial do campo Ey em t = 0
    Ey[0, :] = np.sin(np.pi / a * xgrid)

    # marcha no tempo
    for n in range(nt - 1):

        # atualiza o H
        for i in range(nx - 1):
            Hz[n+1, i] = Hz[n, i] + (dt / (mu0 * dx)) * (Ey[n, i+1] - Ey[n, i])

        # atualiza o E
        for i in range(1, nx - 1):
            # por causa do i-1 no final que o range começa em 1
            Ey[n+1, i] = Ey[n, i] + (dt / (e0 * dx)) * (Hz[n+1, i] - Hz[n+1, i-1])

        # zero na condição de contorno (extremidades)
        Ey[n+1, 0] = 0
        Ey[n+1, -1] = 0


    \end{lstlisting}

    \chapter{Código Exercício 3}\label{cap:codigo_ex3}

    \begin{lstlisting}[language=Python, breaklines=true, basicstyle=\tiny]
    
    import numpy as np
    import matplotlib.pyplot as plt
    import seaborn as sns
    from scipy.fft import fft, fftfreq

    # --- parâmetros físicos ---
    eps = 8.854187817e-12  # Vacuum permittivity (F/m)
    mu = 4 * np.pi * 1e-7  # Vacuum permeability (H/m)
    c = 1 / np.sqrt(eps * mu)

    # --- domínio ---
    a = 0.5
    Nx = Ny = 300
    dx = dy = 5e-3
    cx, cy = Nx//2, Ny//2

    dt = 10e-12
    Nt = 2000

    # --- campos ---
    Ez = np.zeros((Nt, Nx, Ny))
    Hx = np.zeros((Nt, Nx, Ny))
    Hy = np.zeros((Nt, Nx, Ny))

    # mascara para a cavidade
    x = np.arange(Nx) * dx
    y = np.arange(Ny) * dy
    X, Y = np.meshgrid(x, y, indexing='ij')

    # circle parameters
    R = a               # radius (meters)

    # boolean mask: True inside the circle
    mask = (X - cx * dx)**2 + (Y - cy * dy)**2 <= R**2

    # ponta de prova para coletar pontos para a FFT no centro da cavidade
    probe = []

    for n in range(Nt-1):
        # --- fonte (pulso gaussiano no tempo) ---
        t0 = 50
        spread = 15
        source_gaussiano = np.exp(-((n - t0)**2) / spread**2)

        # --- fonte (senoidal no tempo) ---
        frequency = 1e9  # 1 GHz
        source_senoidal = np.sin(2 * np.pi * frequency * n * dt)

        Ez[n, cx, cy] += source_senoidal

        Hx[n+1, :, :-1] = Hx[n, :, :-1] - (dt/(mu*dy)) * (
            Ez[n, :, 1:] - Ez[n, :, :-1]
        )

        Hy[n+1, :-1, :] = Hy[n, :-1, :] + (dt/(mu*dx)) * (
            Ez[n, 1:, :] - Ez[n, :-1, :]
        )

        # --- Update E field (from H at time n+1) ---
        Ez[n+1, 1:-1, 1:-1] = Ez[n, 1:-1, 1:-1] + (dt/eps) * (
            (Hy[n+1, 1:-1, 1:-1] - Hy[n+1, :-2, 1:-1]) / dx -
            (Hx[n+1, 1:-1, 1:-1] - Hx[n+1, 1:-1, :-2]) / dy
        )

        # zero nas extremidades e para fora da cavidade (condição de contorno)
        Ez[n+1][~mask] = 0

        # salva o valor no centro para a FFT
        probe.append(Ez[n+1, cx, cy])

    # 2. Compute the FFT
    yf = fft(probe)
    xf = fftfreq(len(probe), dt)

    # 3. Plot the results (Positive frequencies only)
    plt.plot(xf[:len(probe)//8], 1e3 * 2.0/len(probe) * np.abs(yf[:len(probe)//8]))
    plt.title("Espectro FFT do campo Ez no centro da cavidade")
    plt.xlabel(r"Frequência $\cdot 10^{10}$ Hz")
    plt.ylabel(r"Amplitude $\cdot 10^{-3}$")
    plt.grid()
    plt.show()
    \end{lstlisting}

\end{anexosenv}
